\expandafter\ifx\csname CAPETINCLUDEMODE\endcsname\relax
\documentclass[11pt,a4paper]{article}
% preamble.tex — packages et commandes communes a toutes les fiches CAPET SII IE
\input{glyphtounicode}
\pdfgentounicode=1
\usepackage{cmap}
\usepackage[T1]{fontenc}
\usepackage[utf8]{inputenc}
\usepackage[french]{babel}
\usepackage[a4paper, top=1.6cm, bottom=1.6cm, left=1.4cm, right=1.4cm]{geometry}
\setlength{\headheight}{14.5pt}
\usepackage{amsmath, amssymb, amsfonts}
\usepackage{siunitx}
\sisetup{locale=FR, output-decimal-marker={,}, inter-unit-product=\ensuremath{{}\cdot{}}}
\usepackage{array, booktabs, multirow, tabularx}
\usepackage{xcolor}
\usepackage[most, breakable]{tcolorbox}
\usepackage{multicol}
\usepackage{enumitem}
\usepackage{mathtools}
\usepackage{lmodern}
\usepackage{microtype}
\usepackage{fancyhdr}
\usepackage{tikz}
\usetikzlibrary{positioning,arrows.meta}
\usepackage{lastpage}

% ── Couleurs par domaine ──────────────────────────────────────────────
\definecolor{colElec}{RGB}{26, 82, 148}       % Électrotechnique — bleu foncé
\definecolor{colElecLight}{RGB}{210, 225, 245} % fond bleu clair
\definecolor{colAuto}{RGB}{34, 120, 66}        % Automatique — vert
\definecolor{colMeca}{RGB}{185, 90, 20}        % Mécanique — orange
\definecolor{colTherm}{RGB}{160, 40, 40}       % Thermique — rouge
\definecolor{colQT}{RGB}{232, 244, 233}        % fond question type (vert très léger)
\definecolor{colJury}{RGB}{255, 243, 205}      % fond jury (jaune ambre)
\definecolor{colJuryFrame}{RGB}{200, 130, 0}   % bordure jury
\definecolor{colGrayLight}{RGB}{245, 245, 247} % fond gris neutre
\definecolor{colSeparator}{RGB}{180, 180, 180} % gris séparateur

% ── Style tcolorbox de section principale ────────────────────────────
\tcbset{
  ficheSection/.style={
    breakable,
    enhanced,
    colback=#1!7!white,
    colframe=#1,
    colbacktitle=#1,
    coltitle=white,
    fonttitle=\bfseries\normalsize,
    left=4pt, right=4pt, top=3pt, bottom=4pt,
    before skip=6pt, after skip=6pt,
  }
}

% ── Style encadré question type ──────────────────────────────────────
\tcbset{
  questiontype/.style={
    breakable,
    enhanced,
    colback=colQT,
    colframe=colElec!55!black,
    colbacktitle=colElec!55!black,
    coltitle=white,
    fonttitle=\bfseries\small,
    left=4pt, right=4pt, top=3pt, bottom=3pt,
    before skip=5pt, after skip=5pt,
  }
}

% ── Style encadré jury ───────────────────────────────────────────────
\tcbset{
  jury/.style={
    breakable,
    enhanced,
    colback=colJury,
    colframe=colJuryFrame,
    fonttitle=\bfseries\normalsize,
    left=4pt, right=4pt, top=3pt, bottom=3pt,
  }
}

% ── Mise en page ─────────────────────────────────────────────────────
\pagestyle{fancy}
\fancyhf{}
\renewcommand{\headrulewidth}{0.4pt}
\renewcommand{\footrulewidth}{0.3pt}

\setlist[itemize]{leftmargin=1.2em, itemsep=1pt, topsep=2pt, parsep=0pt}
\setlist[enumerate]{leftmargin=1.2em, itemsep=3pt, topsep=2pt, parsep=0pt}

\setlength{\parskip}{2pt}
\setlength{\columnsep}{10pt}

\def\CAPETBEGINFILE{\begin{document}}
\def\CAPETENDFILE{\end{document}}
\else
\def\CAPETBEGINFILE{}
\def\CAPETENDFILE{}
\fi

\fancyhead[L]{\textbf{\textcolor{colElec}{T04}} - Harmoniques, THD, puissance deformante}
\fancyhead[R]{\small\textit{CAPET SII IE - Session 2026}}
\fancyfoot[C]{\small\thepage\,/\,\pageref{LastPage}}

\CAPETBEGINFILE

% -----------------------------------------------------------------------
%  BLOC EN-TeTE
% -----------------------------------------------------------------------

\begin{tcolorbox}[
  enhanced, breakable,
  colback=colElec,
  colframe=colElec,
  left=6pt, right=6pt, top=5pt, bottom=5pt,
  before skip=0pt, after skip=4pt
]
  {\color{white}\bfseries\large T04 \quad Harmoniques reseau, THD et puissance deformante}\\[2pt]
  {\color{colElecLight}\small
    \textbf{Domaine :} electrotechnique / Qualite reseau \quad
    \textbf{Presence :} \textbf{10/13} sujets (2013-2026) - theme \emph{tres frequent}
  }\\[1pt]
  {\color{colElecLight!80!white}\footnotesize
    Sous-themes : decomposition de Fourier - THD - facteur de puissance - MLI - norme IEC~61000-2-2 - filtre LC
  }
\end{tcolorbox}

% -----------------------------------------------------------------------
%  1. RAPPELS DE COURS
% -----------------------------------------------------------------------

\begin{tcolorbox}[
  ficheSection=colElec,
  title={1. Rappels de cours},
]

% -- 1.1 Decomposition de Fourier -------------------------------------
\textbf{1.1 Decomposition de Fourier d'un signal periodique}

\smallskip
Tout courant periodique se decompose en :
\[
  i(t) = I_0 + \sum_{n=1}^{\infty} I_n\sqrt{2}\cos(n\omega t + \varphi_n)
\]
oe $I_0$ est la composante continue, $I_1$ le fondamental ($n=1$) et $I_n$ les harmoniques de rang $n$.

\smallskip
\begin{center}
\renewcommand{\arraystretch}{1.45}
\begin{tabular}{lll}
\toprule
\textbf{Grandeur} & \textbf{Expression} & \textbf{Remarque} \\
\midrule
Valeur efficace totale & $I_{eff} = \sqrt{I_1^2 + I_2^2 + \cdots} = I_1\sqrt{1+THD^2}$ & racine de la somme des carres \\[5pt]
THD courant (IEC) & $THD = \dfrac{\sqrt{I_2^2+I_3^2+\cdots}}{I_1}$ & normalise par $I_1$, sans unite \\[8pt]
THD en pourcentage & $THD\% = THD \times 100$ & limite norme : $25\%$ \\
\bottomrule
\end{tabular}
\end{center}

\medskip
\hrule\vspace{4pt}

% -- 1.2 Puissances en presence d'harmoniques ------------------------
\textbf{1.2 Puissances en presence d'harmoniques}

\smallskip
Lorsque la tension reseau est sinusoedale ($v(t) = V_1\sqrt{2}\cos\omega t$), seul le fondamental du courant contribue - $P$ et $Q$ :

\begin{center}
\renewcommand{\arraystretch}{1.45}
\begin{tabular}{llll}
\toprule
\textbf{Grandeur} & \textbf{Formule} & \textbf{Unite} & \textbf{Remarque} \\
\midrule
Puissance active & $P = V_1 I_1 \cos\varphi_1$ & W & seul le fondamental contribue \\[4pt]
Puissance reactive & $Q = V_1 I_1 \sin\varphi_1$ & VAR & lie au dephasage fondamental \\[4pt]
Puissance deformante & $D = V_1 \sqrt{I_2^2+I_3^2+\cdots} = V_1 I_1 \cdot THD$ & VAD & due aux harmoniques \\[4pt]
Puissance apparente & $S = \sqrt{P^2+Q^2+D^2}$ & VA & \\[4pt]
Facteur de puissance & $FP = \dfrac{P}{S} = \dfrac{\cos\varphi_1}{\sqrt{1+THD^2}}$ & - & $FP \leq \cos\varphi_1$ \\
\bottomrule
\end{tabular}
\end{center}

\medskip
\hrule\vspace{4pt}

% -- 1.3 MLI et filtrage ----------------------------------------------
\textbf{1.3 MLI et filtrage des harmoniques}

\smallskip
\begin{center}
\renewcommand{\arraystretch}{1.35}
\begin{tabular}{lp{8cm}}
\toprule
\textbf{Mode} & \textbf{Effet sur le spectre harmonique} \\
\midrule
Pleine onde & Harmoniques de rangs $h = 6k\pm1$ (5, 7, 11, 13e) - basse frequence. $THD \approx 48\text{}59\%$. \\[4pt]
MLI ($f_{MLI}$) & Harmoniques de basse frequence elimines ; repousses autour de $f_{MLI}$ et ses multiples. Filtrage LC passe-bas possible avec $f_c < f_{MLI}$. \\[4pt]
Filtre LC & $f_c = \dfrac{1}{2\pi\sqrt{LC}}$ - placer $f_c \ll f_{MLI}$ pour attenuer les harmoniques MLI. \\
\bottomrule
\end{tabular}
\end{center}

\smallskip
{\footnotesize \textbf{Norme IEC/EN~61000-2-2} : $THD < 25\%$ (soit $THD = 0{,}25$) pour le raccordement au reseau public.}

\end{tcolorbox}

% -----------------------------------------------------------------------
%  2. FORMULAIRE SYNTHeTIQUE
% -----------------------------------------------------------------------

\begin{tcolorbox}[
  ficheSection=colElec,
  title={2. Formulaire synthetique},
]
\small
\begin{multicols}{2}
\renewcommand{\arraystretch}{1.45}
\begin{tabular}{@{}p{4.0cm}ll@{}}
\toprule
\textbf{Formule} & \textbf{Usage} & \textbf{U.} \\
\midrule
\multicolumn{3}{@{}l}{\footnotesize\textit{Harmoniques}} \\
$THD = \sqrt{\sum_{n\geq2}I_n^2}\,/\,I_1$ & Taux distorsion & - \\
$I_{eff} = I_1\sqrt{1+THD^2}$ & Valeur efficace & A \\
$THD_{po} \approx 0{,}48$ & Pleine onde mono & - \\
$THD_{triphase} \approx 0{,}48$ & Pont triphase th. & - \\
\midrule
\multicolumn{3}{@{}l}{\footnotesize\textit{Puissances}} \\
$P = V_1 I_1 \cos\varphi_1$ & P active & W \\
$Q = V_1 I_1 \sin\varphi_1$ & P reactive & VAR \\
$D = V_1 I_1 \cdot THD$ & P deformante & VAD \\
$S = \sqrt{P^2+Q^2+D^2}$ & P apparente & VA \\
\bottomrule
\end{tabular}

\columnbreak

\renewcommand{\arraystretch}{1.45}
\begin{tabular}{@{}p{4.0cm}ll@{}}
\toprule
\textbf{Formule} & \textbf{Usage} & \textbf{U.} \\
\midrule
\multicolumn{3}{@{}l}{\footnotesize\textit{Facteur de puissance}} \\
$FP = P/S$ & Definition & - \\
$FP = \cos\varphi_1/\sqrt{1+THD^2}$ & Avec harm. & - \\
\midrule
\multicolumn{3}{@{}l}{\footnotesize\textit{Norme et filtre}} \\
$THD < 0{,}25$ & Limite IEC & - \\
$f_c = 1/(2\pi\sqrt{LC})$ & Coupure filtre & Hz \\
$h = 6k\pm1$ & Rangs harm. tri. & - \\
\midrule
\multicolumn{3}{@{}l}{\footnotesize\textit{Ondulation hacheur}} \\
$\Delta i = E\alpha(1-\alpha)/(Lf)$ & Ondulation DC & A \\
\bottomrule
\end{tabular}
\end{multicols}
\normalsize
\end{tcolorbox}

% -----------------------------------------------------------------------
%  3. QUESTIONS TYPES
% -----------------------------------------------------------------------

\begin{tcolorbox}[
  ficheSection=colElec,
  title={3. Questions types - extraites des sujets reels},
]

% -- QT1 -------------------------------------------------------------

% Mini-procedure de resolution type copie concours
\begin{tcolorbox}[questiontype,title={Procedure type copie concours (5 etapes)}]
\begin{enumerate}[leftmargin=1.5em]
  \item \textbf{Qualifier le modele} : hypotheses, domaine de validite, grandeurs d'entree/sortie.
  \item \textbf{Ecrire les lois} : equations de base avec symboles et unites SI.
  \item \textbf{Isoler l'inconnue} : expression litterale, puis application numerique tracee.
  \item \textbf{Controler l'ordre de grandeur} : coherence physique et unite finale.
  \item \textbf{Conclure concours} : phrase courte orientee decision technique.
\end{enumerate}
\end{tcolorbox}

% Visuel de methode : chemin de resolution CAPET
\begin{center}
\begin{tikzpicture}[font=\small, node distance=7mm and 6mm, >=Stealth]
  \node[draw, rounded corners, fill=colElec!8, inner sep=4pt] (n1) {Notions};
  \node[draw, rounded corners, fill=colElec!8, right=of n1, inner sep=4pt] (n2) {Methode};
  \node[draw, rounded corners, fill=colElec!8, right=of n2, inner sep=4pt] (n3) {Application annale};
  \node[draw, rounded corners, fill=colElec!8, right=of n3, inner sep=4pt] (n4) {Verification jury};
  \draw[->, thick, colElec] (n1) -- (n2);
  \draw[->, thick, colElec] (n2) -- (n3);
  \draw[->, thick, colElec] (n3) -- (n4);
\end{tikzpicture}
\end{center}
\begin{tcolorbox}[questiontype,
  title={QT1 - THD pleine onde vs MLI 6 kHz - conformite norme IEC \hfill\normalfont\footnotesize E2 disciplinaire 2024 (B9-B13)}
]
\textbf{enonce.}
L'onduleur alimente la MAS. En mode pleine onde, le spectre donne $I_5/I_1$, $I_7/I_1$, etc.
En mode MLI $f_{MLI}=\SI{6}{\kilo\hertz}$, le THD mesure est $0{,}25$.
\textbf{B9.} Calculer $THD_\text{cas1}$. \textbf{B10.} Comparer - la norme IEC~61000-2-2. \textbf{B11-B13.} Conclure sur l'effet de la MLI et proposer un filtre.

\tcbline

\textbf{Corrige.}

\textbf{Cas 1 - pleine onde.}
\[
  THD_\text{cas1} = \sqrt{\left(\frac{I_5}{I_1}\right)^2 + \left(\frac{I_7}{I_1}\right)^2 + \cdots} \approx \mathbf{0{,}59} \quad (59\%)
\]
Non conforme : $59\% > 25\%$ (limite IEC~61000-2-2).

\textbf{Cas 2 - MLI \SI{6}{\kilo\hertz}.}
\[
  THD_\text{cas2} = \mathbf{0{,}25} \quad (25\%)
\]
Conforme - la limite. La MLI deplace les harmoniques de basse frequence vers $f \approx f_{MLI} = \SI{6}{\kilo\hertz}$, qui peuvent ensuite etre attenues par un filtre LC passe-bas ($f_c < f_{MLI}$).
\end{tcolorbox}

\vspace{4pt}

% -- QT2 -------------------------------------------------------------

\begin{tcolorbox}[questiontype,
  title={QT2 - Puissance deformante et facteur de puissance reel \hfill\normalfont\footnotesize E2 sujet 2018 (questions 7-9)}
]
\textbf{enonce.}
Un convertisseur absorbe un courant deforme : $I_1$, $\varphi_1$, $THD$ connus. La tension reseau est sinusoedale ($V_1$). Calculer $D$, $S$ et $FP$.

\tcbline

\textbf{Corrige.}

\textbf{etape 1 - puissances elementaires.}
\[
  P = V_1 I_1 \cos\varphi_1
  \qquad
  Q = V_1 I_1 \sin\varphi_1
  \qquad
  D = V_1 \cdot I_1 \cdot THD
\]

\textbf{etape 2 - puissance apparente et facteur de puissance.}
\[
  S = \sqrt{P^2+Q^2+D^2}
  \qquad
  FP = \frac{P}{S} = \frac{\cos\varphi_1}{\sqrt{1+THD^2}}
\]
\textit{Remarque :} $\cos\varphi_1$ est le facteur de \emph{deplacement} (seul dephasage fondamental) ; $FP$ est le facteur de puissance \emph{reel}, toujours inferieur - $\cos\varphi_1$ en presence d'harmoniques.
\end{tcolorbox}

\vspace{4pt}

% -- QT3 -------------------------------------------------------------

\begin{tcolorbox}[questiontype,
  title={QT3 - THD d'un onduleur pleine onde monophase \hfill\normalfont\footnotesize E2 disciplinaire 2025 (CV3-2)}
]
\textbf{enonce.}
L'onduleur CV3-2 fonctionne en pleine onde. Calculer le THD du fondamental de la tension de sortie.

\tcbline

\textbf{Corrige.}

La serie de Fourier d'une tension pleine onde est :
\[
  v(t) = \frac{4V_{DC}}{\pi}\!\left(\cos\omega t - \frac{1}{3}\cos 3\omega t + \frac{1}{5}\cos 5\omega t - \cdots\right)
\]

\textbf{Calcul du THD.}
Les amplitudes harmoniques valent $V_n = V_1/n$, donc :
\[
  THD_{fond} = \frac{\sqrt{(V_1/3)^2+(V_1/5)^2+(V_1/7)^2+\cdots}}{V_1}
             = \sqrt{\frac{1}{9}+\frac{1}{25}+\frac{1}{49}+\cdots} \approx \mathbf{0{,}48}
\]
\end{tcolorbox}

\end{tcolorbox}

% -----------------------------------------------------------------------
%  4. ATTENTES DU JURY
% -----------------------------------------------------------------------

\begin{tcolorbox}[
  jury,
  title={4. Attentes du jury - Erreurs recurrentes et points de vigilance},
]

\begin{enumerate}[label=\textbf{\arabic*.}, leftmargin=*, itemsep=5pt]

  \item \textbf{Normalisation du THD par $I_1$, pas par $I_{eff}$.}
  La definition IEC est $THD = \sqrt{\sum_{n\geq2}I_n^2}\,/\,I_1$.
  Diviser par $I_{eff}$ donne un resultat different - erreur frequente.

  \item \textbf{La MLI ne supprime pas les harmoniques, elle les deplace.}
  Elle repousse les harmoniques vers $f_{MLI}$ et ses multiples. Un filtre LC reste necessaire.

  \item \textbf{$FP \neq \cos\varphi_1$ en presence d'harmoniques.}
  $FP = \cos\varphi_1/\sqrt{1+THD^2}$ - toujours inferieur - $\cos\varphi_1$.
  Confondre les deux fausse toute conclusion sur la conformite.

  \item \textbf{Puissance active et harmoniques.}
  Seul le fondamental contribue - $P$ si la tension est sinusoedale.
  Les harmoniques du courant vont uniquement dans $D$ (puissance deformante).

  \item \textbf{Norme IEC~61000-2-2 : limite $25\%$ sur le THD global.}
  Cette limite s'applique au $THD$ total du courant, pas - chaque harmonique individuelle.

  \item \textbf{Ondulation $\neq$ THD.}
  L'ondulation $\Delta i$ est une grandeur temporelle crete--crete. Le THD est frequentiel. Ne pas les confondre dans les conclusions.

  \item \textbf{Rangs harmoniques pont triphase.}
  $h = 6k\pm1$ (5, 7, 11, 13e). Les rangs 3, 9, 15e sont absents en triphase equilibre.

\end{enumerate}

\medskip
\begin{tcolorbox}[
  colback=colJuryFrame!12!colJury,
  colframe=colJuryFrame,
  left=4pt, right=4pt, top=2pt, bottom=2pt,
  fontupper=\small\itshape
]
Rapport jury 2024 : "La MLI permet de limiter le THD du courant et donc d'en limiter les ondulations. La mise en place de la MLI genere des harmoniques dans les hautes frequences, permettant l'utilisation d'un filtre passe-bas."\\
Rapport jury 2025 : "La partie 2 (convertisseurs et harmoniques) a ete abordee par 97\,\% des candidats mais seulement 8{,}5\,\% ont obtenu la moyenne. Le calcul de THD reste mal maetrise."
\end{tcolorbox}

\end{tcolorbox}

% -----------------------------------------------------------------------
%  PIED - PReSENCE DANS LES SUJETS
% -----------------------------------------------------------------------

\begin{tcolorbox}[
  enhanced, breakable,
  colback=colGrayLight,
  colframe=colSeparator,
  left=4pt, right=4pt, top=3pt, bottom=3pt,
  before skip=4pt, after skip=0pt,
  fonttitle=\small\bfseries,
  title={Presence dans les sujets analyses (2013-2026)},
]
\begin{center}
\footnotesize\renewcommand{\arraystretch}{1.15}
\begin{tabular}{lccccccccccccc}
\toprule
\textbf{Annee} & 13 & 14 & 15 & 16 & 17 & 18 & 19 & 20 & 21 & 22 & 23 & 24 & 25 \\
\midrule
E1 (commune)   & - & - & - & - & - & - & - & - & $\checkmark$ & $\checkmark$ & $\checkmark$ & $\checkmark$ & $\checkmark$ \\
E2 discipl. IE & - & - & - & - & $\checkmark$ & $\checkmark$ & $\checkmark$ & $\checkmark$ & $\checkmark$ & $\checkmark$ & $\checkmark$ & $\checkmark$ & $\checkmark$ \\
\bottomrule
\end{tabular}
\end{center}
\smallskip
\footnotesize\textbf{Sous-themes frequents :}
THD (10/13) - comparaison pleine onde / MLI (6/13) - facteur de puissance avec harmoniques (8/13) - norme IEC~61000-2-2 (5/13) - filtre LC (5/13)

\medskip
\textbf{Ordres de grandeur (reperes rapides).}
\begin{itemize}[leftmargin=1.3em]
  \item Reseau BT FR : tension phase-neutre typique \SI{230}{\volt}, tension phase-phase typique \SI{400}{\volt}, frequence \SI{50}{\hertz}.
  \item Courants d'etude frequents en sujet : de quelques \si{\ampere} a quelques dizaines de \si{\ampere}.
  \item En conversion d'energie, un rendement attendu reste generalement entre 0{,}8 et 0{,}98 selon la technologie.
\end{itemize}

\medskip
\textbf{Synthese express (a refaire sans support).}
\begin{enumerate}[leftmargin=1.4em]
  \item Poser les hypotheses du modele et les conventions de signe.
  \item Ecrire 2 a 4 relations cles avec unites.
  \item Realiser un calcul numerique complet et verifier l'ordre de grandeur.
  \item Justifier le resultat par un argument technique de niveau CAPET.
\end{enumerate}

\end{tcolorbox}

\CAPETENDFILE









